\chapter{Roteiro do vídeo pitch}

\section{Introdução — Problema}
\textbf{Tempo reservado: 2 a 3 minutos}

\begin{itemize}
    \item Abrir com o gancho de dor: em uma fábrica de envase rápido, garrafas com tampa ausente, tampa mal rosqueada ou deformidade no corpo chegam sem triagem até a etapa final da linha.
    \item Explicar a consequência concreta: isso trava a esteira principal, gera paradas não programadas, derramamento de líquido e desperdício de lotes já embalados.
    \item Situar o problema no processo: mostrar (com imagem, diagrama ou vídeo curto) o ponto exato da linha onde a falha ocorre — a etapa final, antes do empacotamento.
    \item Fechar a introdução com a pergunta que a solução vai responder: "como identificar e remover esses itens defeituosos antes que cheguem a essa etapa, sem parar a esteira?"
\end{itemize}

\section{Solução — Proposta tecnológica}
\textbf{Tempo reservado: 4 a 5 minutos}

\begin{itemize}
    \item Apresentar a proposta em uma frase: inspeção visual multi-vista (topo + duas laterais) com trigger físico confiável, classificação automática dos três tipos de defeito e ejeção confirmada por sensor.
    \item Explicar o fluxo entrada → processamento → saída (reaproveitar o diagrama do capítulo de arquitetura): trigger (E18-D80NK + VL53L0X) → captura sincronizada (3 câmeras) → classificação por vista → fusão por votação → registro no hub → ejeção confirmada.
    \item Indicar, sem entrar em detalhe técnico excessivo, como cada componente resolve uma parte específica do problema:
    \begin{itemize}
        \item o trigger garante que toda garrafa seja capturada, mesmo transparente ou com líquido;
        \item as três vistas cobrem tampa e corpo simultaneamente;
        \item a fusão por votação evita decisão baseada em uma única vista com erro;
        \item o atuador só ejeta com confirmação física, evitando remover item bom por engano.
    \end{itemize}
    \item Reforçar que o valor esperado (redução de paradas, menos desperdício) é uma expectativa a ser validada por ensaio, não um resultado já medido.
\end{itemize}

\section{Demonstração — Funcionamento}
\textbf{Tempo reservado: 5 a 6 minutos}

\begin{itemize}
    \item Anunciar o que será mostrado ao vivo: passagem de garrafas golden sample (uma OK e as anomalias conhecidas: sem tampa, tampa mal rosqueada, deformidade) pelo rig montado.
    \item Indicar o espaço físico reservado: bancada com esteira/rig montado, câmeras posicionadas e iluminação estroboscópica .
  \item Descrever a sequência da demonstração:
    \begin{enumerate}
        \item passagem da garrafa golden OK — deve seguir sem ejeção;
        \item passagem das três garrafas com defeito conhecido, uma de cada vez — cada uma deve ser classificada e ejetada;
        \item mostrar em paralelo o dashboard atualizando o status de cada item em tempo de execução.
    \end{enumerate}
    \item Reservar o tempo de forma explícita no roteiro: cerca de 1 minuto por golden sample (4 no total) + 1 a 2 minutos para mostrar o dashboard e o registro gerado.
    \item Caso algum defeito não seja capturado ao vivo (falha de rede, hardware etc.), ter um vídeo de backup gravado previamente como plano B, deixado claro na narração que é uma gravação de apoio.
\end{itemize}

\section{Conclusão — Fechamento}
\textbf{Tempo reservado: 2 minutos}

\begin{itemize}
    \item Retomar o problema em uma frase: paradas não programadas e desperdício causados por garrafas defeituosas na etapa final da linha.
    \item Retomar a solução em uma frase: inspeção multi-vista com classificação automática e ejeção confirmada, sem interromper o fluxo da esteira.
    \item Apresentar o resultado esperado com o cuidado já adotado no documento de requisitos: falar em termos de expectativa e não de dado comprovado, por exemplo: "esperamos que essa abordagem reduza o número de paradas não programadas causadas por esse tipo de defeito, algo que pretendemos validar com ensaios de throughput e resiliência nas próximas etapas."
    \item Indicar continuidade: quais camadas ainda dependem de validação (ex.: camadas condicionais de detecção de anomalia desconhecida, generalização para outros objetos) para deixar claro o que é núcleo pronto e o que é extensão futura.
    \item Fechar com agradecimento e call-to-action apropriado (perguntas, disponibilidade para detalhamento técnico).
\end{itemize}

\section{Distribuição de tempo (referência rápida)}
\begin{tabularx}{\textwidth}{p{0.3\textwidth} Y}
\toprule
Parte & Tempo reservado\\
\midrule
Introdução & 2–3 min\\
Solução & 4–5 min\\
Demonstração & 5–6 min\\
Conclusão & 2 min\\
\midrule
\textbf{Total} & \textbf{13–16 min (ajustar para caber em 15 min)}\\
\bottomrule
\end{tabularx}

